\begin{frame}{Truth Tables}
\begin{center}
\vspace{1.5cm}
{\Huge\textbf{Truth Tables}}
\end{center}
\end{frame}

% --------------------------------------------------
% Question 1
% --------------------------------------------------

\begin{frame}{Question 1}
A software system grants access only when a user enters a valid password
and successfully completes two-factor authentication.

Let:

$$
P = \text{Valid password}
$$

$$
Q = \text{Two-factor authentication successful}
$$

Represent the access condition using a logical expression and construct its
truth table. Identify the conditions under which access is granted.
\end{frame}

\begin{frame}{Question 1 -- Answer}
The system requires both conditions to be satisfied.

Therefore, the logical expression is:

$$
P \land Q
$$

There are four possible combinations of truth values for $P$ and $Q$.
\end{frame}

\begin{frame}{Question 1 -- Answer}
The truth table is:

$$
\begin{matrix}
P & Q & P\land Q \\
T & T & T \\
T & F & F \\
F & T & F \\
F & F & F
\end{matrix}
$$

The expression is true only when both $P$ and $Q$ are true.
\end{frame}

\begin{frame}{Question 1 -- Answer}
Therefore, access is granted only when:

$$
P=T,\qquad Q=T
$$

This means:

\begin{center}
\textit{The password is valid and two-factor authentication is successful.}
\end{center}

For every other combination, access is denied.
\end{frame}

% --------------------------------------------------
% Question 2
% --------------------------------------------------

\begin{frame}{Question 2}
A backup system raises an alert when either the primary storage device
fails or the backup storage device fails.

Let:

$$
P = \text{Primary storage fails}
$$

$$
Q = \text{Backup storage fails}
$$

Represent the alert condition using a logical expression and construct its
truth table. Determine when the alert is not raised.
\end{frame}

\begin{frame}{Question 2 -- Answer}
The alert is raised when at least one of the two storage devices fails.

Therefore, the logical expression is:

$$
P \lor Q
$$

There are four possible combinations of truth values for $P$ and $Q$.
\end{frame}

\begin{frame}{Question 2 -- Answer}
The truth table is:

$$
\begin{matrix}
P & Q & P\lor Q \\
T & T & T \\
T & F & T \\
F & T & T \\
F & F & F
\end{matrix}
$$

The expression is false only when both $P$ and $Q$ are false.
\end{frame}

\begin{frame}{Question 2 -- Answer}
Therefore, the alert is not raised only when:

$$
P=F,\qquad Q=F
$$

This means:

\begin{center}
\textit{Neither the primary nor the backup storage device has failed.}
\end{center}
\end{frame}

% --------------------------------------------------
% Question 3
% --------------------------------------------------

\begin{frame}{Question 3}
An intrusion detection system follows the rule:

\begin{center}
\textit{If suspicious activity is detected, the security team must be notified.}
\end{center}

Let:

$$
P = \text{Suspicious activity is detected}
$$

$$
Q = \text{Security team is notified}
$$

Represent the rule using an implication and construct its truth table.
Identify the condition that represents a violation of the rule.
\end{frame}

\begin{frame}{Question 3 -- Answer}
The rule has the form:

\begin{center}
\textit{If $P$, then $Q$.}
\end{center}

Therefore, the logical expression is:

$$
P \rightarrow Q
$$

\end{frame}

\begin{frame}{Question 3 -- Answer}
The truth table is:

$$
\begin{matrix}
P & Q & P\rightarrow Q \\
T & T & T \\
T & F & F \\
F & T & T \\
F & F & T
\end{matrix}
$$

An implication is false only when the antecedent is true and the
consequent is false.
\end{frame}

\begin{frame}{Question 3 -- Answer}
The rule is violated when:

$$
P=T,\qquad Q=F
$$

Therefore:

\begin{center}
\textit{Suspicious activity is detected, but the security team is not notified.}
\end{center}

This is the only combination that makes the implication false.
\end{frame}

% --------------------------------------------------
% Question 4
% --------------------------------------------------

\begin{frame}{Question 4}
A cloud service considers a server healthy when both its CPU usage is within
the permitted range and its memory usage is within the permitted range.

However, an emergency override can also declare the server healthy.
\end{frame}

\begin{frame}{Question 4}
Let:

$$
P = \text{CPU usage is within the permitted range}
$$

$$
Q = \text{Memory usage is within the permitted range}
$$

$$
R = \text{Emergency override is active}
$$

The health condition is:

$$
(P \land Q) \lor R
$$

Construct the complete truth table and determine how many combinations
result in the server being considered healthy.
\end{frame}

\begin{frame}{Question 4 -- Answer}
The logical expression is:

$$
(P \land Q) \lor R
$$

There are three propositions.

Therefore, the number of possible combinations is:

$$
2^3=8
$$

So, the complete truth table contains eight rows.
\end{frame}


\begin{frame}{Question 4 -- Answer}
The truth table is:

$$
\begin{matrix}
P & Q & R & P\land Q & (P\land Q)\lor R \\
T & T & T & T & T \\
T & T & F & T & T \\
T & F & T & F & T \\
T & F & F & F & F \\
F & T & T & F & T \\
F & T & F & F & F \\
F & F & T & F & T \\
F & F & F & F & F
\end{matrix}
$$

\end{frame}

\begin{frame}{Question 4 -- Answer}
The final column is true in four cases.

Therefore:

$$
\boxed{4\text{ combinations}}
$$

result in the server being considered healthy.
\end{frame}

\begin{frame}{Question 4 -- Answer}
The server is considered healthy when either:

\begin{itemize}
\item both CPU and memory conditions are satisfied, or
\item the emergency override is active.
\end{itemize}

The truth table examines every possible combination of the three
conditions.
\end{frame}

% --------------------------------------------------
% Question 5
% --------------------------------------------------

\begin{frame}{Question 5}
An autonomous vehicle follows a safety policy:

\begin{center}
\textit{If an obstacle is detected, then the vehicle must either brake or
change direction.}
\end{center}

Let:

$$
P = \text{Obstacle is detected}
$$

$$
Q = \text{Vehicle brakes}
$$

$$
R = \text{Vehicle changes direction}
$$

Represent the policy using a logical expression and construct its complete
truth table. Identify all situations in which the safety policy is violated.
\end{frame}

\begin{frame}{Question 5 -- Answer}
The vehicle must brake or change direction whenever an obstacle is detected.

Therefore, the consequent is:

$$
Q \lor R
$$

The complete policy is:

$$
P \rightarrow (Q \lor R)
$$

\end{frame}

\begin{frame}{Question 5 -- Answer}
Since there are three propositions, the truth table contains:

$$
2^3=8
$$

possible combinations.

We first evaluate:

$$
Q\lor R
$$

\end{frame}

\begin{frame}{Question 5 -- Answer}
The complete truth table is:

$$
\begin{matrix}
P & Q & R & Q\lor R & P\rightarrow(Q\lor R) \\
T & T & T & T & T \\
T & T & F & T & T \\
T & F & T & T & T \\
T & F & F & F & F \\
F & T & T & T & T \\
F & T & F & T & T \\
F & F & T & T & T \\
F & F & F & F & T
\end{matrix}
$$

\end{frame}

\begin{frame}{Question 5 -- Answer}
The policy is violated only when:

$$
P=T,\qquad Q=F,\qquad R=F
$$

In this situation:

\begin{center}
\textit{An obstacle is detected, but the vehicle neither brakes nor changes
direction.}
\end{center}
\end{frame}

\begin{frame}{Question 5 -- Answer}
Therefore, the safety policy is violated in exactly one of the eight
possible situations.

$$
\boxed{P=T,\;Q=F,\;R=F}
$$

This shows how a truth table can systematically identify unsafe
combinations of system conditions.
\end{frame}
